\chapter{Metodologia}
\label{cap:metodologia}

% Introdução breve do capítulo explicando que a pesquisa tem caráter empírico e quantitativo, 
% focando na comparação de três padrões de decomposição em três aplicações distintas.

Este capítulo descreve a metodologia experimental adotada para conduzir a avaliação comparativa de estratégias de decomposição de monólitos em microsserviços. A pesquisa caracteriza-se por uma abordagem empírico-quantitativa, fundamentada no projeto, implementação, refatoração e instrumentação de três aplicações monolíticas distintas, submetidas a três metodologias de decomposição arquitetural.

A estrutura metodológica foi desenhada para capturar dados empíricos ao longo do processo de migração, utilizando o padrão \textit{Strangler Fig} como mecanismo unificado de extração incremental. A coleta de dados ocorre de forma iterativa por meio de \textit{snapshots} formais, permitindo avaliar o impacto de cada estratégia sob quatro dimensões: agilidade/DevOps, custo/infraestrutura, desempenho e manutenibilidade.

As seções a seguir detalham a estrutura experimental da pesquisa, a justificativa e caracterização dos estudos de caso, os critérios teóricos para seleção dos padrões de decomposição, os procedimentos operacionais da migração incremental por \textit{snapshots} e a matriz de métricas e instrumentos de medição utilizados.

\section{Desenho da Pesquisa e Abordagem Experimental}
\label{sec:desenho_pesquisa}


\section{Seleção e Justificativa das Aplicações (Estudos de Caso)}
\label{sec:selecao_aplicacoes}
% Justificar a escolha de domínios diferentes para testar os gargalos das arquiteturas.





\subsection{Aplicação Crítica: Sistema de Inventário}
% Focar na necessidade de zero downtime e integridade no apontamento do chão de fábrica.

https://github.com/egp1020/inventory-api


\subsection{Aplicação de Alta Performance: Processamento de Concorrência}
% Focar no volume de requisições e vazão (throughput).

https://github.com/schoibs/auction


\subsection{Aplicação Simples: E-commerce (Catálogo e Carrinho)}
% Focar no CRUD tradicional e na demanda de leitura rápida.

https://github.com/sappChak/nestjs-ecommerce


\section{Procedimentos de Migração e Marcos de Avaliação}
\label{sec:procedimentos_snapshots}
% Descrever como o experimento será conduzido iterativamente.

\subsection{O Padrão \textit{Strangler Fig} como Roteador}
% Como o tráfego será desviado sem desligar o monólito.

\subsection{Definição dos \textit{Snapshots} de Transição}
% Explicar os 4 estágios de medição:
% Snapshot 0 (Baseline/Monólito)
% Snapshot 1 (1ª Extração)
% Snapshot 2 (2ª Extração)
% Snapshot Final (Decomposição Concluída)

